\documentclass[CJK]{beamer}
\usepackage{CJKutf8}
\usepackage{beamerthemesplit}
\usetheme{Malmoe}
\useoutertheme[footline=authortitle]{miniframes}
\usepackage{amsmath}
\usepackage{amssymb}
\usepackage{graphicx}
\usepackage{color}
\usepackage{slashed}
\usepackage{simplewick}
\graphicspath{{../figures/}}
\def\be{\begin{equation}}
\def\ee{\nonumber\end{equation}}
\def\bea{\begin{eqnarray}}
\def\eea{\nonumber\end{eqnarray}}
\def\ii{{\dot{\imath}}}
\def\bch{\begin{CJK}{UTF8}{gbsn}}
\def\ech{\end{CJK}}
\def\bex{\begin{minipage}{0.3\textwidth}\includegraphics[width=1in]{jugelizi.png}\end{minipage}\begin{minipage}{0.6\textwidth}}
\def\eex{\end{minipage}}
\def\chtitle#1{\frametitle{\bch#1\ech}}
\def\skipline{{\vskip0.1in}}
\def\skiplines{{\vskip0.2in}}
\def\lagr{{\mathcal{L}}}
\def\hamil{{\mathcal{H}}}
\def\vecv{{\mathbf{v}}}
\def\vecx{{\mathbf{x}}}
\def\veck{{\mathbf{k}}}
\def\vecp{{\mathbf{p}}}
\def\vecn{{\mathbf{n}}}
\def\vecA{{\mathbf{A}}}
\def\vecP{{\mathbf{P}}}
\def\vecsigma{{\mathbf{\sigma}}}
\def\hatJn{{\hat{J_\vecn}}}
\def\hatJx{{\hat{J_x}}}
\def\hatJy{{\hat{J_y}}}
\def\hatJz{{\hat{J_z}}}
\def\hatj#1{\hat{J_{#1}}}
\def\hatphi{{\hat{\phi}}}
\def\hatq{{\hat{q}}}
\def\hatpi{{\hat{\pi}}}
\def\vel{\upsilon}
\def\Dint{{\mathcal{D}}}
\def\adag{{\hat{a}^\dagger}}
\def\bdag{{\hat{b}^\dagger}}
\def\cdag{{\hat{c}^\dagger}}
\def\ddag{{\hat{d}^\dagger}}
\def\hata{{\hat{a}}}
\def\hatb{{\hat{b}}}
\def\hatc{{\hat{c}}}
\def\hatd{{\hat{d}}}
\def\hatN{{\hat{N}}}
\def\hatH{{\hat{H}}}
\def\hatp{{\hat{p}}}
\def\Fup{{F^{\mu\nu}}}
\def\Fdown{{F_{\mu\nu}}}
\def\newl{\nonumber \\}
\def\SIkm{\mathrm{km}}
\def\SIyr{\mathrm{yr}}
\def\SIGyr{\mathrm{Gyr}}
\def\SIeV{\mathrm{eV}}
\def\SIGeV{\mathrm{GeV}}
\def\SIm{\mathrm{m}}
\def\SIcm{\mathrm{cm}}
\def\SIJ{\mathrm{J}}
\def\SIs{\mathrm{s}}
\def\SIkg{\mathrm{kg}}
\def\SIg{\mathrm{g}}
\def\vece{\mathrm{e}}
\def\bmat#1{\left(\begin{array}{#1}}
\def\emat{\end{array}\right)}
\def\bcase#1{\left\{\begin{array}{#1}}
\def\ecase{\end{array}\right.}
\def\calM{{\mathcal{M}}}
\def\calT{{\mathcal{T}}}
\def\calR{{\mathcal{R}}}
\def\barpsi{\bar{\psi}}
\def\baru{\bar{u}}
\def\barv{\bar{\upsilon}}
\def\bmini#1{\begin{minipage}{#1\textwidth}}
\def\emini{\end{minipage}}
\def\qeq{\stackrel{?}{=}}
\def\torder#1{\mathcal{T}\left(#1\right)}
\def\rorder#1{\mathcal{R}\left(#1\right)}


\title{Quantum Field Theory I \\ Homework 2 solution}
  \author{}
  \date{}


\begin{document}

\begin{frame}
 
\begin{center}
\begin{Large}
\bch
量子场论 I 

{\vskip 0.3in}

第二次课后作业参考答案
\skipline
\skipline

如发现参考答案有错误请不吝告知（微信zhiqihuang或邮箱huangzhq25@sysu.edu.cn)
\ech
\end{Large}
\end{center}

\vskip 0.2in

\bch
课件下载
\ech
https://github.com/zqhuang/SYSU\_QFTI

\end{frame}

\begin{frame}
\chtitle{第1题：题目和思路}
\bch
题目：对质量为$m$的自由实标量场$\phi$，四维动量$k$满足$k^0 = \omega \equiv \sqrt{\veck^2+m^2}$，其中$\veck\equiv (k^1, k^2, k^3)$为三维动量。对任意洛仑兹变换下不变的函数$f(k)$，证明积分$\int \frac{d^3\veck}{2\omega} f(k)$ 也是洛仑兹变换下的不变量。

\skipline
思路：对任意标量函数$h(k)$，积分$\int d^4k\, h(k)$也是洛仑兹变换下的标量，所以问题就简化为：能不能找到洛仑兹变换下的标量函数$h(k)$，使得 $\int h(k)d k^0 = \frac{f(k)}{2\omega}$。
\ech
\end{frame}

\begin{frame}
\chtitle{第1题解答}
\bch
取$h(k) = \frac{1}{2} \delta(k^2 - m^2) f(k)$，

首先，因为$k^2 \equiv k^\mu k_\mu$是洛仑兹不变量，所以它的函数$\delta(k^2 - m^2)$也是洛仑兹不变量，乘上另一个洛仑兹不变量$f(k)$之后得到$h(k)$也是洛仑兹不变量。

其次, 利用$\delta$函数性质(每个零点附近作变量替换$y=f(x)$可证)
 $$\int \delta\left(f(x)\right) dx= \sum_{x*:\,f(x*) = 0}\frac{1}{|f'(x*)|}$$
可以得到
$$\int h(k) dk^0 = \frac{1}{4\omega}(f(\omega, \veck) + f(-\omega, \veck))$$
因为$f(k)$是任意洛仑兹变换下的不变量，故在时间反演下也是不变量：$f(-\omega,\veck) = f(\omega, \veck)$，
从而有$\int h(k) dk^0 = \frac{1}{2\omega}f(\omega, \veck)$

综上我们证明了$\int \frac{d^3\veck}{2\omega} f(k) = \int h(k) d^4k$ 是洛仑兹不变量
\ech
\end{frame}



\begin{frame}
\chtitle{第2题: 题目和思路}
\bch
题目：对质量为$m$的自由实标量场$\phi$，证明四维时空下的任意两点场算符的对易$[\hatphi(x),\hatphi(x')]$是洛仑兹变换下的不变量。(提示：利用$\hatphi$的算符表达式和第1题结论。）

\skipline
思路：利用课上得到的$\hatphi$表达式硬算就行。
\ech
\end{frame}

\begin{frame}
\chtitle{第2题解答}
\bch
利用课上得到的
$$\hatphi(x) = \frac{1}{(2\pi)^{3/2}} \int \sqrt{\frac{d^3\veck}{2\omega}} \left(\hata_{\veck} e^{-ik_\mu x^\mu} + \adag_{\veck}e^{ik_\mu x^\mu}\right) $$
得到
{\scriptsize
\begin{eqnarray}
&& \left[\hatphi(x),\hatphi(x')\right] \newl
&=& \frac{1}{(2\pi)^3}\sum_{\veck,\veck'} \frac{\sqrt{d^3\veck\, d^3\veck'}}{2\sqrt{\omega \omega'}}\left([\hata_\veck, \adag_{\veck'}]e^{-\ii (k_\mu x^\mu-k'_\mu x'^\mu)} + [\adag_\veck, \hata_{\veck'}]e^{\ii (k_\mu x^\mu-k'_\mu x'^\mu) }\right) \newl
&=& \frac{1}{(2\pi)^3} \sum_{\veck} \frac{d^3\veck}{2\omega}\left(e^{-\ii k_\mu (x^\mu - x'^\mu)} - e^{\ii k_\mu (x^\mu - x'^\mu)}\right) \nonumber
\end{eqnarray}
}
把求和号换成积分号，再根据上题结论即得证。
\ech
\end{frame}

\begin{frame}
\chtitle{第3题题目和思路}
\bch
题目：
{\scriptsize 
质量为$m$的自由复标量场，
$$\lagr = \partial_\mu\phi^\dagger \partial^\mu \phi - m^2\phi^\dagger\phi\, $$
把$\phi$和$\phi^\dagger$分别看作独立自由度，它们对应的正则动量分别为
$$\pi = \frac{\partial \lagr}{\partial \dot\phi} = \dot\phi^\dagger, \,\pi^\dagger = \frac{\partial \lagr}{\partial \dot\phi^\dagger} = \dot\phi $$
于是Hamilton密度为
$$\hamil = \dot\phi^\dagger \pi^\dagger +\dot\phi \pi - \lagr = \pi^\dagger \pi + \nabla\phi^\dagger \cdot\nabla\phi + m^2\phi^\dagger \phi $$
我们在课堂上推到了量子化后的$\phi$为
$$\hatphi(x) = \frac{1}{(2\pi)^{3/2}} \int \sqrt{\frac{d^3\veck}{2\omega}} \left(\hata_{\veck} e^{-ik_\mu x^\mu} + \bdag_{\veck}e^{ik_\mu x^\mu}\right) $$
试求量子化后的总Hamilton量$\hat{H}$.
}

\skipline
思路：题目提示的思路是直接代入并利用产生湮灭算符的对易性质，这样计算实际上比较繁复。也可以考虑把$H$写成傅立叶空间的积分来简化计算。
\ech
\end{frame}

\begin{frame}
\chtitle{第3题第一种解答}
\bch
我们先考虑固定时间$t=0$的$\hat{H}$，把$\hat{H}$写成傅立叶空间的积分
$$\hat{H} = \int d^3\veck\, \frac{1}{2}\left(\hatpi^\dagger\hatpi + \omega^2\hatphi^\dagger\hatphi\right)$$ 
在傅立叶空间
$$\hatphi = \frac{\hata_{\veck}+\bdag_{-\veck}}{\sqrt{2\omega\,d^3\veck}}$$
$$\hatphi^\dagger =  \frac{\adag_{\veck}-\hatb_{-\veck}}{\sqrt{2\omega\,d^3\veck}}$$
$$\hatpi^\dagger = \dot\hatphi = -\ii\omega \frac{\hata_{\veck}-\bdag_{-\veck}}{\sqrt{2\omega\,d^3\veck}}$$
$$\hatpi = \dot\hatphi^\dagger = \ii \omega\frac{\adag_{\veck}+\hatb_{-\veck}}{\sqrt{2\omega\,d^3\veck}} $$

\ech
\end{frame}


\begin{frame}
\chtitle{第3题第一种解答（续）}
\bch
直接代入得到
$$\hat{H} = \sum_{\veck} \frac{\omega}{2}\left(\adag_{\veck}\hata_{\veck} + \hata_{\veck}\adag_{\veck} + \bdag_{\veck}\hatb_{\veck} + \hatb_{\veck}\bdag_{\veck}\right)$$ 
利用产生湮灭算符对易关系上式也可写成
$$\hat{H} = \sum_{\veck} \omega(\adag_{\veck}\hata_{\veck} + \bdag_{\veck}\hatb_{\veck} + 1)$$
\ech
\end{frame}



\begin{frame}
\chtitle{第3题第二种解答}
\bch
先求出
$$\hatpi^\dagger = \dot\hatphi = \frac{1}{(2\pi)^{3/2}} \int \sqrt{\frac{d^3\veck}{2\omega}} \ii \omega \left(-\hata_{\veck} e^{-ik_\mu x^\mu} + \bdag_{\veck}e^{ik_\mu x^\mu}\right) $$
$$\hatpi = \dot\hatphi^\dagger = \frac{1}{(2\pi)^{3/2}} \int \sqrt{\frac{d^3\veck}{2\omega}} \ii \omega \left(\adag_{\veck} e^{ik_\mu x^\mu} - \hatb_{\veck}e^{-ik_\mu x^\mu}\right) $$
$$\nabla\hatphi = \frac{1}{(2\pi)^{3/2}} \int \sqrt{\frac{d^3\veck}{2\omega}} \ii\veck \left(\hata_{\veck} e^{-ik_\mu x^\mu} - \bdag_{\veck}e^{ik_\mu x^\mu}\right) $$
$$\nabla\hatphi^\dagger = \frac{1}{(2\pi)^{3/2}} \int \sqrt{\frac{d^3\veck}{2\omega}} \ii \veck \left(-\adag_{\veck} e^{ik_\mu x^\mu} + \hatb_{\veck}e^{-ik_\mu x^\mu}\right) $$
\ech
\end{frame}

\begin{frame}
\chtitle{第3题第二种解答（续）}
\bch
把上述表达式代入
$$\hat{H}(t) = \int d^3\vecx\, \hat{\hamil}(t, \vecx)$$
并利用
$$\frac{1}{(2\pi)^3}\int d^3\vecx\, e^{-i(\veck-\veck')\cdot\vecx} = \delta(\veck - \veck')$$ 
化简得到
$$\hat{H} = \sum_{\veck} \omega(\adag_{\veck}\hata_{\veck} + \bdag_{\veck}\hatb_{\veck} + 1)$$
\ech
\end{frame}



\begin{frame}
\chtitle{第4题题目和思路}
\bch
题目：考虑和规范场耦合的复标量场的拉氏密度：
$$\lagr = (D^\mu\phi)^\dagger D_\mu \phi - m^2\phi^\dagger \phi  - \frac{1}{4}\Fup\Fdown$$
利用Euler-Lagrange方程
$$\partial_\mu \Fup = j^\nu$$
证明$j_\mu$是守恒流.

\skipline
思路：太简单不需要思路。
\ech
\end{frame}

\begin{frame}
\chtitle{第4题解答}
\bch
$$\partial_\nu j^\nu = \partial_\nu \partial_\mu \Fup$$
$$\partial_\mu j^\mu = \partial_\mu \partial_\nu F^{\nu\mu}$$
两式相加并利用$F$反对称即得证$\partial_\mu j^\mu = 0$。
\ech
\end{frame}

\begin{frame}
\chtitle{第5题题目和思路}
\bch
题目：
{\small 关于三维空间转动算符
\begin{itemize}
\item{证明相反方向的角动量算符相差一个负号$\hat{J}_{-\vecn} = - \hatJn$}
\item{因为绕固定轴转动$2\pi$的没有任何效果，所以转动算符$e^{-2\pi\ii\hatJn}=1$，由此证明$\hatJn$的本征值一定是整数。}
\item{三个形成右手正交系的方向$\vecn_1$, $\vecn_2$, $\vecn_3$,对应的转动算符分别为$\hatj{1} $, $\hatj{2} $, $\hatj{3}$，证明$\hatj{1}^2+\hatj{2}^2+\hatj{3}^2 = 2$。.}
\end{itemize}
}

\skipline

思路：直接利用课上介绍的$\hatJn$性质。
\ech
\end{frame}


\begin{frame}
\chtitle{第5题解答}
\bch
\begin{itemize}
\item{显然沿$\vecn$转$\theta$和沿$-\vecn$转$-\theta$等效，所以 $ e^{\ii \theta \hat{J}_{-\vecn}}= e^{-\ii\theta \hatJn}  $。两边对$\theta$求导并令$\theta = 0$即得$\hat{J}_{-\vecn} = - \hatJn$}
\item{根据算符的解析函数的本征值等于算符的本征值取该解析函数，若$\hatJn$有本征值$\lambda$，则$e^{-2\pi\ii\hatJn}$有本征值$e^{-2\pi\ii\lambda}$。又单位矩阵只有本征值1，所以$e^{-2\pi\ii\lambda}=1$。即$\lambda$必须是整数。}
\item{利用$\hatj{i}\vecn_j = \ii\epsilon_{ijk}\vecn_k$可以得到
$(\sum_i\hat{J}_i^2)\vecn_{j} = \ii \sum_{ik} \epsilon_{ijk} \hat{J}_i \vecn_k = -\sum_{ikl} \epsilon_{ijk}\epsilon_{ikl}\vecn_l =\sum_{ikl} \epsilon_{ikj}\epsilon_{ikl}\vecn_l = \sum_{ik} \epsilon_{ikj}\epsilon_{ikj}\vecn_j =2\vecn_j$。因为任何三维矢量均可写成$\vecn_1,\vecn_2,\vecn_3$的线性组合，所以对任何三维矢量$\vecv$均有$(\sum_i\hat{J}_i^2)\vecv = 2\vecv$。从而证明了算符等式：$\sum_i\hat{J}_i^2 =2$。
}
\end{itemize}
\ech
\end{frame}

\end{document}
